\chapter{Conclusions}

\section{Summary}

This work set out to design and implement a real-time air quality monitoring pipeline capable of continuously ingesting sensor readings from 414 Ukrainian cities, detecting pollution events and delivering alerts with sub-minute latency, providing live dashboard visualisation, and persisting both raw readings and aggregated results for historical analysis. All four objectives stated in Section~\ref{sec:objectives} were achieved.

The pipeline is organised around five loosely coupled stages. A FastAPI-based data producer publishes sensor readings to Apache Kafka, which provides durable, replayable buffering between production and consumption. An Apache Flink job reads from Kafka and performs two concurrent tasks: windowed aggregation over 60-second tumbling windows and Z-score-based anomaly detection with a ten-window onset and recovery state machine. Aggregated results are persisted to TimescaleDB, which serves as the data source for a Grafana dashboard providing real-time visualisation of conditions across all monitored stations. Raw readings are archived to Amazon S3 via the Kafka Connect S3 Sink connector, forming an immutable data lake for future reprocessing. A lightweight alerting consumer delivers onset and recovery notifications by email when the stream processor detects sustained anomalies.

The entire stack is reproducible with a single Docker Compose command, with startup ordering enforced through health checks and sensitive configuration passed through environment variables.

\section{Evaluation Results}

The pipeline was evaluated through a series of load tests using a historical dataset replayed at configurable speed. End-to-end latency was measured as the elapsed time between window close and database insertion — the delay between a pollution event occurring and its data becoming queryable.

At production-realistic replay rates (5× through 1,000×, corresponding to 2,070 to 414,000 records per minute), p95 latency stabilised at approximately 5 to 9 seconds, well within the sub-minute threshold. The latency floor corresponds to the configured 5-second watermark delay plus Flink processing and database write time. At 1× speed, sparse inter-event intervals caused watermark stalls that raised p95 to 37 seconds; this is an artifact of the single-reading-per-minute regime rather than a limitation of the pipeline at realistic load. At 10,000× the data producer itself became the bottleneck, capping effective throughput at approximately 1.79 million records per minute; even at this rate p95 latency remained at 13 seconds, well within the one-minute requirement.

Alert detection operated correctly across all tests: onset notifications were dispatched after ten consecutive elevated windows and recovery notifications after ten consecutive normal ones, confirming that the state machine suppresses spurious alerts for sustained events without introducing false negatives.
